% ch3_b §3.3 (书页 83–90 / p097–p104)
%--A-TAIL--
%JOIN%格点处。它对应于电子集中分布在原子上的状态。但这正是我们所预期的。$\mathcal{V}_{G}$ 的\emph{负}值对应于在每个原子邻域内为负的周期势——一个\emph{吸引}电子的势。这样，$\psi^{-}$ 由于与原子的势阱相配合，能量被降低；类似地，$\psi^{+}$ 的能量被升高，因为它使电子密度在正势区域变大。仿照原子轨道的称谓，我们说 $\psi^{-}$ 是 $s$ 型的，$\psi^{+}$ 是 $p$ 型的。

顺便指出，展开式 (3.12) 使 $\psi_{\mathbf{k}}$ 能够满足 Bloch 定理 (1.58)。我们可以写成
\begin{align*}
\psi_{\mathbf{k}} &= e^{i\mathbf{k}\cdot\mathbf{r}}\sum_{\mathbf{g}}\alpha_{\mathbf{k}-\mathbf{g}}e^{-i\mathbf{g}\cdot\mathbf{r}}\notag \\
&= e^{i\mathbf{k}\cdot\mathbf{r}}u_{\mathbf{k}}(\mathbf{r}), \tag{3.23}
\end{align*}
其中 $u_{\mathbf{k}}(\mathbf{r})$ 是晶格中的周期函数，正如 (1.60) 中那样。

%FIG{40}: {周期势，以及 $s$ 型与 $p$ 型态中的电子密度。}

我们还注意到，图 39 中 $A$ 与 $A'$ 两处的波函数是完全相同的；二者均由 (3.21) 给出。类似地，$\psi^{+}(\frac{1}{2}G)$ 与 $\psi^{+}(-\frac{1}{2}G)$ 除相差一个常数外也相同，因此 $B$ 与 $B'$ 两点在图 38 与图 39 的区图式中是等价的。这是 §1.5 中所讨论的那些原理的一个实例，在那里已经证明：一个态的波矢量在相差一个倒格矢的意义下是任意的。

\section{近自由电子模型}

一维情形的这些原理可以推广到三维。如果我们沿着 k 空间中某个特定方向考察 $\mathcal{E}(\mathbf{k})$ 的行为，就会发现它沿着一条自由电子抛物线变化，而在穿过布里渊区边界时发生跳跃式的不连续（图 41）。每个不连续的大小将取决于势中相应 Fourier 分量的强弱。如果区边界是倒格矢 $\mathbf{G}$ 的垂直平分面，那么不连续的大小将是 $2|\mathcal{V}_{\mathbf{G}}|$ 的量级。每当一条等能线穿过区边界时，同样会出现这种不连续。于是（图 42），本应穿过 $O$ 点的自由电子球被分裂成双曲面的一叶——在区内 $S$ 处，由 (3.17) 的一个根生成——以及边界之外另一支的一叶 $S'$。

%FIG{41}: {沿 k 空间中的 $OAB$ 线，在区边界处 $\mathcal{E}(\mathbf{k})$ 出现不连续。}

%FIG{42}: {区边界处等能线的不连续。}

这意味着函数 $\mathcal{E}(\mathbf{k})$ 在第一布里渊区内部是连续的——但在任何一点越过区边界时，能量都有一个跳跃。这使该理论最引人注目的推论之一成为可能：\emph{对某些能量值，可能不存在电子态}。能谱可以出现能隙。电子能级分成了若干\emph{能带}。这是一条基本性质，它衍生出金属和半导体的许多性质。它是 Bragg 衍射的一个后果。电子波函数被“拉”得与晶格具有相同的周期性，并在能量上分裂开来。

目前我们关心的是求出 $\mathcal{E}(\mathbf{k})$ 的问题。上一节的计算提示了如下的方法：——用方程组 (3.14) 求展开式 (3.12) 中系数的解；把所有与能被 $\mathbf{k}$ 附近的区边界所混合的波的系数相对应的项分离出来，对能量 $\mathcal{E}(\mathbf{k})$ 求解久期行列式（secular determinant）；再用形如 (3.9) 的微扰展开修正势的其他 Fourier 分量的效应。

这称为\emph{近自由电子}（nearly-free-electron，N.F.E.）方法。要使它有效，Fourier 分量 $\mathcal{V}_{\mathbf{g}}$ 的级数应当迅速收敛。可是乍看起来这一点似乎不大可能。$\mathcal{V}(\mathbf{r})$ 是离子阵列的势。我们知道，离子核心附近的场非常强——$\mathcal{V}(\mathbf{r})$ 在每个格点处都有一个向下延伸的长而尖的尖峰。这意味着它具有波长非常短的 Fourier 分量，因此当 $\mathbf{g}$ 为第一区尺度的许多倍时，$\mathcal{V}_{\mathbf{g}}$ 仍可能很大。或者我们也可以说，$\psi_{\mathbf{k}}(\mathbf{r})$ 在离子核心内部必须表现得像原子轨道一样，带有对应于高动能的多次振荡。这同样意味着，在表示式 (3.12) 中我们需要许多项，一直用到非常短的波长。

正是这些论证曾使这一方法作为能带结构的实用计算方案而遭到忽视。但现在看来，通过引入赝势（pseudo-potential）的观念，可以使它在形式上有效。这将由 §3.6 展开。

N.F.E. 模型作为说明函数 $\mathcal{E}(\mathbf{k})$ 在 k 空间中行为的手段非常有用。考虑一个矩形点阵——二维情形——它的倒格子也是矩形的，如图 43 所示。第一布里渊区很容易构造——就是矩形 $RR'R''R'''$。$\mathcal{E}(\mathbf{k})$ 在此区内部必须连续，而且在勾画这个晶胞的任何一条线上都可能出现不连续。

但再考虑中央区之外 $\mathcal{E}(\mathbf{k})$ 不连续点的轨迹。它们可以是倒格子中任一矢量的垂直平分线。因此，图中任何一条线上都可能有不连续——例如 $\mathbf{g}_{3}$ 的垂直平分线就是一条穿过区角隅的对角线。我们最初那张关于不连续的图（图 41）是不对的——它没有画出可能发生的一切，即便是方格子情形也不例外。因此，在图 43 中的 $X_{1},X_{2},X_{3},X_{4}$ 等点处，$\mathcal{E}(\mathbf{k})$ 都可能有跳跃和能隙。

%FIG{43}: {在矩形倒格子中，$\mathcal{E}(\mathbf{k})$ 的不连续将出现在 $X_{1}$、$X_{2}$、$X_{3}$、$X_{4}$ 处。}

这看起来似乎表明 $\mathcal{E}(\mathbf{k})$ 是一个远比实际情形复杂的函数。例如，假定我们试着画出由 N.F.E. 模型得出的一些等能线。它们看上去将如图 44 那样（把注意力限定在这些平面组以内的区域）。我们可以先画出自由电子圆，再按图 42 所示的方式加入不连续，从而把它们构造出来。或者更确切地说，我们应当对该点处相互混合的那些波的能量求解久期行列式。于是在 $P$ 点有与 $P'$ 处的波的混合，因为在自由电子方案中这两点能量相同，且相差一个倒格矢。类似地，为求 $Q$ 点的解，我们需要考虑这样一个 $3\times 3$ 行列式，它描述波矢本应为 $OQ$、$OQ'$ 和 $OQ''$ 的那些波（在自由电子模型中）的混合。

现在的意思是，在 $P$ 与 $P'$ 两点我们将求解的是完全相同的方程。两点的能级相同。我们将其中的较低者归属于第一区内部；在每种情形下，较高的能级都属于区边界之外的等能线。同样，在 $Q$、$Q'$ 和 $Q''$ 处方程也是同一个，并给出三个能量根。我们可以把其中的最低者，在这些点的每一处，归属于图 44 中各边界以内的区域。于是这三个点能量全部相同。我们可以考察这些点的邻域，并类似地证明：$A$ 与 $A'$、$B$ 与 $B'$、$C$ 与 $C'$ 都是在此意义下等价的点对；$A$ 与 $A'$ 是同一波函数的两套可互换的标记，借助倒格子平移即可相互还原。

%FIG{44}: {矩形区中的等能线。}

现在看一看，当我们用倒格矢把第一区之外的那些碎块零料平移过去时会发生什么。它们完美地拼合起来，构成单一的区——在 Brillouin 扩展区图式中，这些部分本来都被说成属于\emph{第二区}。而且，边界上的等价点恰好匹配，使等能线连续起来。\emph{简约区中的能量是 $\mathbf{k}$ 的连续函数}。

我们表示能面的方案于是就像图 46 那样，每条能带配一个区。我们可以把这些等能线想成不同曲面的等高线；把它们一层层堆叠起来，就可以把 $\mathcal{E}(\mathbf{k})$ 表示成简约区中 $\mathbf{k}$ 的多值函数，而不必像扩展区图式中那样让 $\mathcal{E}(\mathbf{k})$ 成为单值函数。

我们还可以再进一步。可以证明（我们的 N.F.E. 模型就能提供一个证明），在每条能带之内 $\mathcal{E}(\mathbf{k})$ 是连续的，且导数也连续。图 45 中画出的那些内部边界在能面上并不显现；等能线平滑地从它们上面越过。现在让我们回想，基本区本身就是倒格子的晶胞，通过倒格子平移可以把它重复排列，产生布满全空间的图案。而且，像 $P$ 与 $P'$ 这样原本等价、因而代表同一状态和相同能量的点，现在相互重合，于是 $\mathcal{E}(\mathbf{k})$ 跨越这些边界也是连续的。可以证明，等能线从一个晶胞到下一个晶胞平滑地衔接。于是，\emph{给定能带中的能量是 $\mathbf{k}$ 的具有倒格子周期的连续函数}。我们把它称为重复区图式（repeated zone scheme）。这与我们在 §2.2 中就恒定晶格频率曲面所观察到的原理如出一辙。

%FIG{45}: {约化成单一的区。}

%FIG{46}: {对应于图 44 的简约区图式。}

这一原理的一个有趣应用，是 Harrison 的方法：在重复区图式中构造价为 $Z$ 的金属的费米面。设 N.F.E. 方案中的微扰势非常小。此时能量面必定是球面。我们算出体积恰为一个区的 $\frac{1}{2}Z$ 倍的球面的半径，并以原点为中心把它画出。再围绕倒格子的每一点画出同样的球面，于是便得到一幅具有重复区图式周期性的图案（图 48）。

现在，我们可以从这幅图案中挑出各种能够连续拼合的部分，组成在每个区中都重复的曲面。这些分立的图形，每一个便构成费米面的一个分支，即费米面在第 1、第 2、第 3……区中的部分（图 49）。当作球面画出时，这些曲面在各部分相接处看似带有尖点，但与势的 Fourier 分量相联系的态的混合会把棱角磨圆，给出光滑的几何对象。正如我们将在第 9 章看到的那样，所得的曲面往往与实验上可观察到的曲面相当相像，而且对于研究复杂金属中的费米面，它们肯定提供了宝贵的指引。

%FIG{47}: {对应于图 46 的重复第一区。}

%FIG{48}: {自由电子费米面的 Harrison 作图。}

这些代数公式和几何构图也出现在动力学衍射理论（dynamical theory of diffraction）中（§2.6）。在非常接近衍射条件 (3.11) 时，穿过晶格的快电子束或 X 射线束就不再能表示成带单一波矢 $\mathbf{k}$ 的简单平面波。各种衍射波必须像 (3.12) 那样包括进波函数之中，从而改变了
%（句子未完，下接书页 91 / p105 页首 "excitation."——应续作“……从而改变了激发的速度”，由下一部分的 B-TAIL 续译。）

%FIG{49}: {对应于图 48 的简约区：第一区（满）、第二区、第三区、第四区。}

%FIG{50}: {X 射线或快电子动力学衍射的色散面。}
